\chapter{Implementación}

\section{Organización del repositorio}
\pending{Relacionar la estructura del repositorio con las capas de la plataforma y separar código, pruebas, experimentos, resultados y documentación.}

\section{Implementación del núcleo experimental}
\pending{Describir configuraciones, validaciones, datasets, entrenamiento, métricas y persistencia implementados.}

\section{Implementación de la QCNN y arquitecturas}
\pending{Explicar la implementación concreta de codificaciones, convolución, pooling lógico, observables, parameter-shift, simulador rápido vectorizado (\texttt{src/fast\_qcnn.py}) y arquitecturas alternativas (\texttt{src/qcnn\_architectures.py}: \texttt{v1}, \texttt{v1-param-pool} y \texttt{qibo}).}

\section{Entrenamiento variacional y ruta QGD}
\pending{Documentar la evolución desde el legado híbrido (PyTorch/Adam) hacia la línea variacional sin PyTorch (\texttt{src/variational\_training.py}: GD, Adam, SPSA y secuencial tipo Rotosolve/NFT) y el prototipo QGD M1--M5 (\texttt{src/quantum\_update.py}, \texttt{src/quantum\_training.py}), detallando la desviación modular reversible y el control clásico remanente.}

\section{Implementación del pipeline de datos}
\pending{Describir particionado estratificado, normalización min-max, PCA en Breast Cancer ajustado solo en train y transformación de etiquetas \texttt{-1/+1}.}

\section{Implementación de la ejecución cuántica y preparación BSC}
\pending{Documentar simulación exacta, shots, ruido local exacto en matriz de densidad frente al modo legado de Qibo 0.1.11, mitigación de lectura, transpilador propio (\texttt{src/transpilation.py}), mock hardware y adaptador \texttt{src/qililab\_adapter.py} (no validado aún en QPU).}

\section{Implementación del modelo y controlador de QML Workbench}
\pending{Describir las capas libres de Qt en \texttt{src/workbench/model/} y \texttt{src/workbench/controller/}, la ejecución CLI sin interfaz (\texttt{src/workbench/run.py}) y la paridad exacta con la fase C.}

\section{Implementación de la vista en PySide6}
\pending{Describir los componentes gráficos de \texttt{src/workbench/view/} (lienzo, paleta, formularios dinámicos, panel de resultados, exportación JSON/CSV y \texttt{GuiDispatcher}).}

\section{Integración entre capas}
\pending{Explicar el flujo completo desde el grafo de QML Workbench y los scripts de \texttt{experiments/} hasta los motores de ejecución y la persistencia de resultados.}

\section{Pruebas}
\pending{Documentar la suite completa de pruebas unitarias y de integración (\texttt{tests/}, incluida la vista Qt en modo \texttt{offscreen}), equivalencia con Qibo, reproducibilidad, cono de luz y preparación BSC.}

\section{Limitaciones de implementación}
\pending{Registrar características no implementadas y decisiones técnicas pendientes.}
